\chapter{From Sound to Vectors}\label{ch:sound}
\chapterlede{Everything in this project begins by turning pressure waves into
arrays of numbers, and arrays of numbers into frequencies. This chapter builds
the discrete Fourier transform from scratch, proves the identities the rest of
the book leans on, and ends at the two workhorse measurements the visualizer
makes sixty times per second: the spectrogram and the autocorrelation pitch
estimate.}

\section{Sampling}

A microphone measures a continuous function $x_c : \R \to \R$ (air pressure
against time). Computers store countably many numbers, so we sample.

\begin{definition}[Sampling]\label{def:sampling}
Given a continuous signal $x_c$ and a \emph{sampling rate} $f_s > 0$ (in Hz),
the sampled signal is the sequence
\[
x[n] = x_c\!\left(\tfrac{n}{f_s}\right), \qquad n \in \Z .
\]
\end{definition}

Sampling destroys information in general: infinitely many continuous signals
agree on any given sample sequence. The classical theorem says that
\emph{bandlimited} signals are the exception.

\begin{theorem}[Nyquist--Shannon]\label{thm:nyquist}
Let $x_c$ be bandlimited to $B$ Hz, meaning its continuous Fourier transform
vanishes outside $[-B, B]$. If $f_s > 2B$, then $x_c$ is uniquely determined
by its samples $x[n]$, and can be reconstructed as
\[
x_c(t) \;=\; \sum_{n \in \Z} x[n] \,
\operatorname{sinc}\!\big(f_s t - n\big),
\qquad \operatorname{sinc}(u) = \frac{\sin(\pi u)}{\pi u}.
\]
\end{theorem}

\begin{proof}[Proof idea]
Bandlimitedness means the Fourier transform $\hat{x}_c$ is supported on
$[-B,B]$. Sampling in time corresponds to \emph{periodizing} in frequency:
the transform of the sample train is $\sum_k \hat{x}_c(f - k f_s)$, a sum of
copies of $\hat{x}_c$ spaced $f_s$ apart. If $f_s > 2B$ the copies do not
overlap, so multiplying by the indicator of $[-f_s/2, f_s/2]$ recovers
$\hat{x}_c$ exactly; that multiplication in frequency is convolution with the
sinc kernel in time.
\end{proof}

When $f_s \le 2B$ the copies overlap and distinct frequencies become
indistinguishable after sampling --- \emph{aliasing}: a tone at frequency $f$
is captured identically to one at $f_s - f$.

\begin{incode}
The project stores two versions of every stem
(\texttt{ingest.py}): a listening copy at $f_s = 44{,}100$ Hz stereo
(Nyquist frequency $22.05$ kHz, above human hearing) and an analysis copy at
$f_s = 22{,}050$ Hz mono (\texttt{config.py}: \texttt{SR = 22050}), which is
cheaper to process and still covers content up to $11$ kHz --- plenty for
beat, pitch, and chroma analysis. Every resample (e.g.\ to MERT's $24$ kHz in
\texttt{embeddings.py}, or CLAP's $48$ kHz in \texttt{backfill\_clap.py})
must low-pass first for exactly the reason in
Theorem~\ref{thm:nyquist}; \texttt{librosa.resample} does this internally.
\end{incode}

\section{The Discrete Fourier Transform}

Fix a frame length $N$ and view a chunk of signal as a vector
$x = (x[0], \dots, x[N-1]) \in \C^N$.

\begin{definition}[DFT]\label{def:dft}
Let $\omega_N = e^{-2\pi i / N}$. The \emph{discrete Fourier transform} of
$x \in \C^N$ is $X \in \C^N$ with
\[
X[k] \;=\; \sum_{n=0}^{N-1} x[n]\, \omega_N^{kn}
\;=\; \sum_{n=0}^{N-1} x[n]\, e^{-2\pi i k n / N},
\qquad k = 0, \dots, N-1 .
\]
\end{definition}

The DFT is the change of basis into the family of sampled complex sinusoids
$e_k[n] = e^{2\pi i k n / N}$. That family is orthogonal:

\begin{lemma}[Orthogonality]\label{lem:orth}
For $k, \ell \in \{0, \dots, N-1\}$,
\[
\sum_{n=0}^{N-1} e^{2\pi i (k-\ell) n / N} =
\begin{cases}
N & k = \ell,\\
0 & k \ne \ell.
\end{cases}
\]
\end{lemma}

\begin{proof}
If $k = \ell$ every term is $1$. Otherwise let
$q = e^{2\pi i (k-\ell)/N} \ne 1$; the sum is geometric:
$\sum_{n=0}^{N-1} q^n = \frac{q^N - 1}{q - 1}$,
and $q^N = e^{2\pi i (k - \ell)} = 1$, so the numerator vanishes.
\end{proof}

\begin{corollary}[Inversion and Parseval]\label{cor:parseval}
$x[n] = \frac{1}{N}\sum_{k=0}^{N-1} X[k]\, e^{2\pi i k n/N}$, and
\[
\sum_{n=0}^{N-1} \abs{x[n]}^2 \;=\; \frac{1}{N} \sum_{k=0}^{N-1}
\abs{X[k]}^2 .
\]
\end{corollary}

\begin{proof}
Substitute Definition~\ref{def:dft} into the claimed inverse and swap the two
finite sums; Lemma~\ref{lem:orth} collapses the inner sum to $N\,
\mathbf{1}[n = m]$. Parseval follows by expanding
$\abs{X[k]}^2 = X[k]\overline{X[k]}$ and applying the lemma once more.
\end{proof}

Parseval is the statement that \emph{energy is preserved}: what the ear hears
as loudness can be measured either per-sample or per-frequency-bin. The
visualizer uses both sides of this identity (RMS level on the time side, band
energy on the frequency side).

\begin{remark}[What a bin means]\label{rem:bins}
Bin $k$ of an $N$-point DFT at sampling rate $f_s$ represents frequency
\[
f_k = \frac{k \, f_s}{N} \quad \text{Hz}, \qquad k \le N/2,
\]
so the frequency resolution is $\Delta f = f_s / N$: longer frames see finer
pitch differences. For a real signal, $X[N-k] = \overline{X[k]}$, so only the
first $N/2 + 1$ bins carry information.
\end{remark}

\begin{incode}
The chroma computation in \texttt{scene.js}
(\texttt{computeChroma}) inverts Remark~\ref{rem:bins}: it walks the
analyser's bins and maps each to its frequency via
\texttt{fr = bin * sr / nfft} before folding onto pitch classes
(Appendix~\ref{ch:temperament}). With \texttt{fftSize = 1024} at
$f_s = 44.1$ kHz, $\Delta f \approx 43$ Hz --- coarse at the bottom of the
piano (the gap between $A1$ and $B\flat1$ is only $3.3$ Hz) and luxurious at
the top. This asymmetry between linear bins and logarithmic musical pitch
recurs throughout the project.
\end{incode}

\section{The Fast Fourier Transform}

Computing Definition~\ref{def:dft} literally costs $N$ multiplications for
each of $N$ bins: $\Theta(N^2)$. The FFT gets the same numbers in
$\Theta(N \log N)$.

\begin{theorem}[Cooley--Tukey]\label{thm:fft}
Let $N$ be even, and split $x$ into even- and odd-indexed halves with DFTs
$E$ and $O$ (each of length $N/2$). Then for $k = 0, \dots, N/2 - 1$:
\[
X[k] = E[k] + \omega_N^{k}\, O[k], \qquad
X[k + N/2] = E[k] - \omega_N^{k}\, O[k].
\]
Consequently, for $N = 2^m$, the DFT can be computed in
$\Theta(N \log N)$ operations.
\end{theorem}

\begin{proof}
Split the defining sum by parity of $n$, writing $n = 2r$ and $n = 2r + 1$:
\[
X[k] = \sum_{r=0}^{N/2-1} x[2r]\, \omega_N^{2rk}
     + \omega_N^{k} \sum_{r=0}^{N/2-1} x[2r+1]\, \omega_N^{2rk}
     = E[k] + \omega_N^k O[k],
\]
using $\omega_N^2 = \omega_{N/2}$. The second identity follows from
$\omega_N^{k + N/2} = -\omega_N^{k}$. The cost recurrence
$T(N) = 2\,T(N/2) + \Theta(N)$ solves to $\Theta(N \log N)$ by the master
theorem.
\end{proof}

This is why every frame size in the project is a power of two
(\texttt{fftSize} of $1024$, $2048$): the browser's
\texttt{AnalyserNode} and NumPy's FFT both rely on this factorization.

\section{Frames, Windows, and the Spectrogram}

Music is not stationary --- the whole point is that its frequency content
changes. The fix is to transform short \emph{frames} and accept a tradeoff.

\begin{definition}[STFT and spectrogram]\label{def:stft}
Fix a frame length $N$, a hop $H \le N$, and a window
$w : \{0,\dots,N-1\} \to [0,1]$. The \emph{short-time Fourier transform} of
$x$ is
\[
X[m, k] = \sum_{n=0}^{N-1} x[mH + n]\, w[n]\, e^{-2\pi i k n / N},
\]
and the \emph{spectrogram} is the matrix of squared magnitudes
$S[m,k] = \abs{X[m,k]}^2$.
\end{definition}

Without a window ($w \equiv 1$), chopping the signal introduces artificial
discontinuities at frame edges, and Definition~\ref{def:dft} smears each true
frequency across all bins (\emph{spectral leakage}). A smooth taper such as
the Hann window
\[
w[n] = \tfrac{1}{2}\Big(1 - \cos\tfrac{2\pi n}{N-1}\Big)
\]
concentrates that leakage into nearby bins at the price of slightly widening
each spectral peak.

\begin{remark}[Time--frequency tradeoff]\label{rem:uncertainty}
$\Delta f = f_s / N$ improves with longer frames while time localization
$\Delta t = N / f_s$ worsens \emph{in exact inverse proportion}:
$\Delta t \cdot \Delta f = 1$ regardless of $N$. This is the discrete shadow
of the uncertainty principle. There is no window that beats it; one can only
choose where to spend it. Musical consequence: a system tuned to see bass
notes precisely (long frames) necessarily blurs drum onsets, and vice versa.
\end{remark}

\begin{incode}
The microscope's spectrogram view (\texttt{microscope.py},
\texttt{spectrogram\_png}) renders $\log(1 + S)$ rather than $S$ --- the log
compresses the enormous dynamic range of music ($\sim$60 dB between a hi-hat
tail and a kick) into the few dozen brightness levels an eye can
distinguish, then maps it through the magma colormap. The scene mode's live
analysis instead uses the browser's \texttt{AnalyserNode}
(\texttt{app.js}, \texttt{bands()}), which computes
Definition~\ref{def:stft} with a built-in Blackman window at
\texttt{fftSize = 1024}: $\Delta t \approx 23$ ms, $\Delta f \approx 43$ Hz
--- the ``spend it on time'' end of Remark~\ref{rem:uncertainty}, correct for
a visualizer chasing drum hits.
\end{incode}

\section{Pitch by Autocorrelation}

The spectrogram spreads a note's energy across its harmonics; for
\emph{monophonic} sources (one note at a time --- bass, a single voice) a
time-domain method finds the fundamental more directly.

\begin{definition}[Autocorrelation]\label{def:acf}
The (unnormalized, windowed) autocorrelation of $x$ at lag $\tau \ge 0$ over
a window of length $W$ is
\[
r[\tau] = \sum_{n=0}^{W-1} x[n]\, x[n + \tau].
\]
\end{definition}

\begin{proposition}[Periodicity peaks]\label{prop:acfpeak}
Suppose $x$ is $P$-periodic on the relevant range, i.e.\ $x[n + P] = x[n]$.
Then $r[\tau] \le r[0]$ for all $\tau$, with equality at every multiple of
$P$: $r[0] = r[P] = r[2P] = \cdots$
\end{proposition}

\begin{proof}
Equality at multiples of $P$ is immediate from periodicity. For the bound,
Cauchy--Schwarz gives
\[
r[\tau] = \ip*{\,x_{0:W},\, x_{\tau:\tau+W}}
\le \norm{x_{0:W}} \, \norm{x_{\tau:\tau+W}} = r[0],
\]
where the final equality uses periodicity again (both windows contain the
same multiset of values when $W$ is a whole number of periods).
\end{proof}

So the algorithm is: compute $r[\tau]$ over a plausible range of lags, find
the peak, and report $f_0 = f_s / \tau^\ast$. Two practical failure modes
matter. First, $r$ peaks equally at $2P, 3P, \dots$, so noise can make the
estimator jump down an octave (choosing $2P$) --- the classic \emph{octave
error}. Second, on silence or noise the global maximum is meaningless, so an
energy gate must precede the search.

\begin{incode}
\texttt{Transport.pitchHz} in \texttt{app.js} is
Definition~\ref{def:acf} verbatim: window $W = 512$ samples, lag range
$f_s/320 \le \tau \le f_s/50$ (i.e.\ candidate fundamentals between $50$ and
$320$ Hz --- the playable range of a bass guitar), and an RMS gate of
$0.012$ before searching. The stutter you once saw in the bass fret marker
was octave error plus frame-to-frame noise; the mitigations in
\texttt{scene.js} (sampling every third frame, level-gating, exponential
smoothing $0.7/0.3$) treat the symptoms. The planned cure --- an offline
pitch tracker with Viterbi smoothing across the whole song --- replaces the
per-frame $\argmax$ with a globally optimal path and will live in the ingest
pipeline.
\end{incode}

\section*{Exercises}

\begin{exercise}
A pure $9$ kHz tone is sampled at $f_s = 16$ kHz with no anti-aliasing
filter. What frequency is heard on playback? Generalize: give the perceived
frequency as a function of true frequency $f$ and $f_s$.
\end{exercise}

\begin{exercise}
Prove the conjugate-symmetry claim of Remark~\ref{rem:bins}: if
$x[n] \in \R$ for all $n$, then $X[N - k] = \overline{X[k]}$. How many real
numbers are needed to store the DFT of a real signal, and why does this not
contradict the DFT being invertible?
\end{exercise}

\begin{exercise}
Complete the proof of Corollary~\ref{cor:parseval} in full detail, justifying
the interchange of summations.
\end{exercise}

\begin{exercise}
Show that the Hann window satisfies
$\sum_{m \in \Z} w[n - mH] = \text{const}$ when $H = N/2$ (perfect overlap-add
at 50\% hop). Why does this property matter for reconstructing audio from a
modified STFT --- for instance, in mask-based source separation
(Chapter~\ref{ch:separation})?
\end{exercise}

\begin{exercise}[computational]
Record or load a bass stem from your library. Implement
Definition~\ref{def:acf} and plot $r[\tau]$ for one frame during a held note.
Identify the fundamental's peak and at least one octave-error peak. Then
break the estimator on purpose: find a frame where the $\argmax$ lands at
$2P$.
\end{exercise}

\begin{exercise}
Using Remark~\ref{rem:uncertainty}: the microscope wants to distinguish
$E1$ ($41.2$ Hz) from $F1$ ($43.7$ Hz) in the bass register. What minimum
frame length (in samples at $22{,}050$ Hz, and in milliseconds) does a plain
DFT need? Roughly how many kick drum hits at $170$ BPM occur within one such
frame? Conclude something about why the project analyzes bass pitch and drum
onsets with different tools.
\end{exercise}
